% AY2022 制御工学 第1問〔I〕（転記）
% 出典: 03_制御工学/original/03_制御工学/2022/2022se.pdf
% 要約: 壁-バネ k-質量 m-ダンパ d-自由端 の系。x1(t) は入力変位（自由端）,
%       x2(t) は質量の変位。運動方程式・伝達関数、ステップ応答・最大値、減衰正弦入力応答。

\section*{第1問〔I〕}

下図の系に関する次の問い (1)〜(4) に答えよ. なお, 質量 $m$, 粘性係数 $d$, バネ定数 $k$
とし, $x_1(t)$ は系への入力となる変位, $x_2(t)$ は質量の変位を表す.
また, 初期状態において系は静止しており, ダッシュポットの動作速度に比例した粘性抵抗が
生じるものとする.

\begin{center}
\begin{tikzpicture}[>=Latex]
  \draw[thick] (0,0.3) -- (0,1.7);
  \foreach \y in {0.4,0.55,...,1.65} \draw[thick] (0,\y) -- (-0.22,\y-0.22);
  \draw[thick,decorate,decoration={zigzag,segment length=7,amplitude=4}] (0,1) -- (1.5,1);
  \node at (0.75,1.5) {$k$};
  \draw[thick,fill=gray!12] (1.5,0.45) rectangle (2.25,1.55);
  \node at (1.875,1.95) {$m$};
  \draw[->,thick] (1.875,0.2) -- (2.5,0.2);
  \node at (2.6,0.0) {$x_2(t)$};
  \draw[thick] (2.25,1) -- (2.6,1);
  \draw[thick] (2.6,0.78) rectangle (3.05,1.22);
  \draw[thick] (2.85,1) -- (3.6,1);
  \node at (2.83,1.95) {$d$};
  \draw[thick] (3.6,1) -- (4.1,1);
  \draw[thick,fill=white] (4.1,1) circle (0.06);
  \draw[->,thick] (4.1,1) -- (4.85,1);
  \node at (4.7,0.7) {$x_1(t)$};
\end{tikzpicture}
\end{center}

\begin{enumerate}[label=(\arabic*)]
  \item 系の運動方程式を求めたうえで, 入力を $x_1(t)$, 出力を $x_2(t)$ として伝達関数を求めよ.

  \item 各パラメータを $m=1$, $d=3$, $k=2$, 初期値をすべて $0$ とし, 入力として $x_1(t)$ に
        ステップ入力を与えたときの出力 $x_2(t)$ を求めよ.

  \item 問い (2) で求めた出力 $x_2(t)$ の最大値とその発生時間を求めよ.

  \item 各パラメータを $m=1$, $d=2$, $k=2$, 初期値をすべて $0$ とし,
        入力 $x_1(t)=e^{-t}\sin(t)$ を加えた（ただし, $t<0$ のとき $x_1(t)=0$）.
        このときの出力 $x_2(t)$ を求め, 充分な時間が経過した後の値を求めよ.
\end{enumerate}
